\subsection{Trace of integral operators with continuous kernel}

In this section, we keep the conventions of the preceding section with Y = X and \(\nu = \mu\). We assume that X is a locally compact topological space that is countable at infinity (TG, I, p. 68, def. 5).

In particular, we identify the spaces \(\mathrm{L}^{2}(\mathrm{X}\times\mathrm{X})\) and \(\mathrm{L}^{2}(\mathrm{X})\widehat{\otimes}_{2}\mathrm{L}^{2}(\mathrm{X})\) (lemma 10 of IV, p. 172). We shall denote by \(\widetilde{f}\) the class in \(\mathrm{L}^{2}(\mathrm{X})\) (resp. in \(\mathrm{L}^{2}(\mathrm{X}\times\mathrm{X})\) ) of a function \(f\in\mathcal{L}^{2}(\mathrm{X})\) (resp. in \(\mathcal{L}^{2}(\mathrm{X}\times\mathrm{X})\) ).

PROPOSITION 20. --- Let \((f_n)_{n \in \mathbf{N}}\) be a sequence of measurable maps from X into a metrizable space F. Suppose that the limit of \(f_n(x)\) exists outside a \(\mu\)-negligible subset of X. There exists a sequence \((\mathrm{C}_m)_{m \in \mathbf{N}}\) of compact subsets of X whose union \(\widetilde{\mathrm{X}}\) satisfies \(|\mu|(\mathrm{X} - \widetilde{\mathrm{X}}) = 0\) , such that the functions \(f_n\) are continuous on \(\mathrm{C}_m\) for all \(n \in \mathbf{N}\) and every \(m \in \mathbf{N}\) , and the sequence \((f_n)_{n \in \mathbf{N}}\) converges uniformly to \(f\) on \(\mathrm{C}_m\) for all \(m \in \mathbf{N}\) .

It follows from th. 2 of INT, IV, p. 175, § 5, no. 4 and from prop. 12, b) of INT, IV, p. 188, § 5, no. 8, that the set of compact subsets C of X such that the functions \(f_n\) are continuous on C for every \(n\) , and such that \((f_n)\) converges uniformly to \(f\) on C, is \(\mu\) -dense in X (INT, IV, p. 189, § 5, no. 8, def. 6). The assertion then follows from the remark of INT, IV, p. 189, § 5, no. 8.

Let \(N\) be a function belonging to \(\mathcal{L}^2 (\mathrm{X}\times \mathrm{X})\) and \(u_{\mathrm{N}}\) the Hilbert--Schmidt map from \(\mathrm{L}^2 (\mathrm{X})\) into \(\mathrm{L}^2 (\mathrm{X})\) with kernel \(N\) (Prop. 19 of IV, p. 173). Since the map \(u_{\mathrm{N}}\) is compact, there exist, by Prop. 9 of IV, p. 156, an element \(M\) of \(\overline{\mathbf{N}}\) and, denoting by \(I\) the set of integers \(n\in \mathbf{N}\) such that \(n\leqslant \mathrm{M}\), orthonormal families \((f_i)_{i\in \mathrm{I}}\) and \((g_{i})_{i\in \mathrm{I}}\) in \(\mathcal{L}^2 (\mathrm{X})\) and a family \((\alpha_{i})_{i\in \mathrm{I}}\) in \(\mathbf{R}_+^*\) such that

\[
u _ {\mathrm{N}} (f) = \sum_ {i \in \mathrm{I}} \alpha_ {i} \langle \widetilde {f} _ {i} | f \rangle \widetilde {g} _ {i}\tag{13}
\]

for all \(f \in \mathrm{L}^2(\mathrm{X})\), where the series converges in \(\mathrm{L}^2(\mathrm{X})\) .

For every \(i \in I\) , we denote by \(h_i \in \mathcal{L}^2(\mathrm{X} \times \mathrm{X})\) the function defined by \(h_i(x, y) = \overline{f_i(x)} g_i(y)\) , so that \(\widetilde{h}_i\) is the class of \(\overline{f}_i \otimes g_i\). By Remark 9 of IV, p. 173, the series with general term \(\alpha_i h_i\) converges to N in \(\mathrm{L}^2(\mathrm{X} \times \mathrm{X})\) .

In the rest of this issue, we assume that N is a continuous function.

PROPOSITION 21. --- Suppose that \(u_{\mathrm{N}}\) has finite trace. There then exists a set \(\widetilde{\mathrm{X}} \subset \mathrm{X}\) whose complement \(\mathrm{X} - \widetilde{\mathrm{X}}\) is \(\mu\) -negligible and a function \(\mathrm{H} \in \mathcal{L}^2(\mathrm{X} \times \mathrm{X})\) satisfying the following conditions:

\begin{enumerate}
\def\labelenumi{(\roman{enumi})}
\item
  For every \((x,y)\in\widetilde{\mathrm{X}}\times\widetilde{\mathrm{X}}\) the family \((\alpha_{i}\overline{f_{i}(x)}g_{i}(y))_{i\in\mathrm{I}}\) is summable in \(\mathbf{C}\) and its sum is \(\mathrm{N}(x,y)\) ;
\item
  For every finite subset J of I and every \((x,y)\in\widetilde{\mathbf{X}}\times\widetilde{\mathbf{X}}\) , we have
\end{enumerate}

\[
\left| \sum_ {i \in \mathrm{J}} \alpha_ {i} h _ {i} (x, y) \right| \leqslant \mathrm{H} (x, y);\tag{14}
\]

\begin{enumerate}
\def\labelenumi{(\roman{enumi})}
\setcounter{enumi}{2}
\tightlist
\item
  The function \(x \mapsto \mathrm{H}(x, x)\) belongs to \(\mathcal{L}^1(\mathrm{X})\) .
\end{enumerate}

Since \(u_{N}\) has finite trace, the family \((\alpha_{i})\) is summable (cor. 4 of IV, p. 165). The series

\[
\sum_ {i \in \mathrm{I}} \alpha_ {i} | f _ {i} | ^ {2}, \quad \sum_ {i \in \mathrm{I}} \alpha_ {i} | g _ {i} | ^ {2}
\]

therefore converge in \(\mathcal{L}^{1}(X)\) and consequently \(\mu\)-almost everywhere. Let F and G denote, respectively, the functions defined \(\mu\)-almost everywhere by the sum of these series, setting \(F(x) = 0\) and \(G(x) = 0\) for all \(x\) such that the corresponding series does not converge. Since F and G belong to \(\mathcal{L}^{1}(X)\), the function H defined by \(H(x,y) = \sqrt{F(x)G(y)}\) belongs to \(\mathcal{L}^{2}(X \times X)\) (INT, V, p. 95, § 8, no. 3, cor. 2).

Let us apply Prop. 20 to the space X, with \(F = \mathbf{C}^{2}\) and with the measurable maps

\[
s_{n}\colon x\mapsto \Bigl (\sum_{\substack{i\in \mathrm{I}\\ i\leqslant n}}\alpha_{i}|f_{i}(x)|^{2},\sum_{\substack{i\in \mathrm{I}\\ i\leqslant n}}\alpha_{i}|g_{i}(x)|^{2}\Bigr)
\]

defined on X. There therefore exists a sequence \((\mathrm{C}_{m})_{m\in\mathbf{N}}\) of compact subsets of X satisfying the following conditions:

\begin{enumerate}
\def\labelenumi{(\arabic{enumi})}
\item
  the union \(\widetilde{X}\) satisfies \(\mu(\mathrm{X}-\widetilde{\mathrm{X}})=0\) ;
\item
  for all \(m \in \mathbf{N}\) and every \(n \in \mathbf{N}\) , the functions \(s_n\) are continuous on \(\mathrm{C}_m\) ;
\item
  for all \(m \in \mathbf{N}\) , the series \(\sum \alpha_{i}|f_{i}|^{2}\) and \(\sum \alpha_{i}|g_{i}|^{2}\) converge uniformly on \(C_{m}\) to F and G, respectively;
\item
  the support of the measure induced by \(\mu\) on \(C_{m}\) is equal to \(C_{m}\) (possibly after replacing \(C_{m}\) by the support of this measure).
\end{enumerate}

Let us prove that the set \(\widetilde{\mathbf{X}}\) and the function H satisfy conditions (i), (ii), and (iii).

Let \((x,y)\in \widetilde{\mathbf{X}}\times \widetilde{\mathbf{X}}\). For every finite subset \(\mathrm{J}\subset \mathrm{I}\) , we have

\[
\left| \sum_ {i \in \mathrm{J}} \alpha_ {i} \overline {{f _ {i} (x)}} g _ {i} (y) \right| \leqslant \left(\sum_ {i \in \mathrm{J}} \alpha_ {i} | f _ {i} (x) | ^ {2}\right) ^ {1 / 2} \left(\sum_ {i \in \mathrm{J}} \alpha_ {i} | g _ {i} (y) | ^ {2}\right) ^ {1 / 2}\tag{15}
\]

whence also

\[
\left| \sum_ {i \in \mathrm{J}} \alpha_ {i} \overline {{f _ {i} (x)}} g _ {i} (y) \right| \leqslant \mathrm{H} (x, y),\tag{16}
\]

which already establishes property (ii).

By inequality (15), the series \(\sum_{i} \alpha_{i} h_{i}\) converges uniformly on \(\mathrm{K}_{m} \times \mathrm{K}_{n}\) for all \((m, n) \in \mathbf{N}^{2}\). Let \(\widetilde{\mathrm{N}}\) be the function on \(\mathrm{X} \times \mathrm{X}\) defined by

\[
\widetilde {\mathrm{N}} (x, y) = \sum_ {i \in \mathrm{I}} \alpha_ {i} h _ {i} (x, y) = \sum_ {i \in \mathrm{I}} \alpha_ {i} \overline {{f _ {i} (x)}} g _ {i} (y)
\]

for all \((x,y)\in\widetilde{\mathrm{X}}\times\widetilde{\mathrm{X}}\) and by \(\widetilde{\mathrm{N}}(x,y)=0\) otherwise. The function \(\widetilde{N}\) is measurable (INT, IV, p. 175, § 5, n°4, th. 2), and it is continuous on \(K_{m}\times K_{n}\) for all \((m,n)\in\mathbf{N}^{2}\) .

From (16), we have \(|\widetilde{\mathrm{N}}(x,y)| \leqslant \mathrm{H}(x,y)\) for all \((x,y) \in \mathrm{X} \times \mathrm{X}\) , and in particular \(\widetilde{\mathrm{N}}\) belongs to \(\mathcal{L}^2(\mathrm{X} \times \mathrm{X})\) (INT, IV, p. 84, § 5, no. 6, th. 5).

Let us prove that \(N = \widetilde{N}\) on \(\widetilde{X} \times \widetilde{X}\) , which will establish property (i). Let f and g be elements of \(\mathcal{L}^{2}(X)\) . We have

\[
\langle g \mid u _ {\widetilde {N}} (f) \rangle = \int_ {X \times X} \widetilde {N} (x, y) f (x) \overline {{g (y)}} d (\mu \otimes \mu) (x, y)
\]

(formula (12) of IV, p. 172). For every \((x,y)\in \widetilde{\mathrm{X}}\times \widetilde{\mathrm{X}}\) and every finite subset J of I, we have

\[
\left| \sum_ {i \in J} \alpha_ {i} \overline {{f _ {i} (x)}} g _ {i} (y) f (x) \overline {{g (y)}} \right| \leqslant | f (x) g (y) | H (x, y)
\]

by formula (16). Since the right-hand side of this inequality is integrable on X × X (INT, V, p. 95, § 8, n° 3, cor. 2), one can apply Lebesgue's theorem (INT, IV, p. 137, § 3, n° 7, th. 6) and formula (13) to deduce that

\[
\begin{array}{c} \langle g | u _ {\widetilde {N}} (f) \rangle = \sum_ {i \in I} \alpha_ {i} \int_ {X \times X} \overline {{f _ {i} (x)}} g _ {i} (y) f (x) \overline {{g (y)}} d (\mu \otimes \mu) (x, y) \\ = \sum_ {i \in I} \alpha_ {i} \langle f _ {i} | f \rangle \langle g | g _ {i} \rangle = \langle g | u _ {N} (f) \rangle . \end{array}
\]

Therefore, it follows that \(u_{\widetilde{\mathrm{N}}} = u_{\mathrm{N}}\) , whence \(\mathrm{N} = \widetilde{\mathrm{N}}\) in \(\mathrm{L}^2(\mathrm{X} \times \mathrm{X})\) (prop. 3 of III, p. 28, b)).

For every \((m,n)\in\mathbf{N}^{2}\), the functions N and \(\widetilde{N}\) are continuous on \(K_{m}\times K_{n}\), therefore are equal on \(K_{m}\times K_{n}\) (prop. 9 of INT, III, p. 69, § 2, n°2) since the support of the measure induced by \(\mu\otimes\mu\) onto \(K_{m}\times K_{n}\) is equal to \(K_{m}\times K_{n}\). Therefore N coincides with \(\widetilde{N}\) on \(\widetilde{X}\times\widetilde{X}\) .

Finally, the upper bound \(\sqrt{\mathrm{FG}} \leqslant (\mathrm{F} + \mathrm{G}) / 2\) implies that the function \(x \mapsto \sqrt{\mathrm{F}(x)\mathrm{G}(x)} = \mathrm{H}(x,x)\) is integrable on X, whence property (iii).

In the remainder of this issue, we keep the notation of the proposition.

Theorem 2. --- Suppose that \(u_{\mathrm{N}}\) has finite trace. Then the function \(x \mapsto \mathrm{N}(x, x)\) belongs to \(\mathcal{L}^1(\mathrm{X})\) and we have

\[
\mathrm{Tr} (u _ {\mathrm{N}}) = \int_ {\mathrm{X}} \mathrm{N} (x, x) d \mu (x).
\]

By conditions (i) and (ii), we have \(|\mathrm{N}(x,x)| \leqslant \mathrm{H}(x,x)\) for \(x \in \widetilde{X}\) , hence the function \(x \mapsto \mathrm{N}(x,x)\) belongs to \(\mathcal{L}^{1}(\mathrm{X})\) by condition (iii).

Condition (i) proves the equality.

\[
\mathrm{N} (x, x) = \sum_ {i \in \mathrm{I}} \alpha_ {i} \overline {{f _ {i} (x)}} g _ {i} (x)
\]

for all \(x \in \widetilde{\mathrm{X}}\) . By conditions (ii) and (iii), one may apply Lebesgue's theorem (INT, IV, p. 137, § 3, no. 7, th. 6), from which it follows that

\[
\begin{array}{r} \int_ {\mathrm{X}} \mathrm{N} (x, x) d \mu (x) = \sum_ {i \in \mathrm{I}} \alpha_ {i} \int_ {\mathrm{X}} \overline {{f _ {i} (x)}} g _ {i} (x) d \mu (x) \\ = \sum_ {i \in \mathrm{I}} \alpha_ {i} \langle f _ {i} | g _ {i} \rangle = \mathrm{Tr} (u _ {\mathrm{N}}), \end{array}
\]

the last equality resulting from the corollary of Proposition 17 of IV, p. 170.

Lemma 11. --- If the endomorphism \(u_{\mathrm{N}}\) of \(\mathrm{L}^2(\mathrm{X})\) is positive, then \(\mathrm{N}(x,x) \geqslant 0\) for all \(x\) in the support of \(\mu\) .

Since \(u_{N}\) is positive, the families \((f_{i})\) and \((g_{i})\) can be chosen such that \(f_{i} = g_{i}\) for all \(i \in I\) (for remark 3 of IV, p. 151). For every \(x \in \widetilde{X}\) , we then have

\[
\mathrm{N} (x, x) = \sum_ {i \in \mathrm{I}} \alpha_ {i} | f _ {i} (x) | ^ {2} \geqslant 0.
\]

Since \(N\) is continuous and \(\mu(\mathbf{X}-\widetilde{\mathbf{X}})=0\), the function \(x \mapsto N(x,x)\) is nonnegative on the support of \(\mu\) .

Remark. --- The converse assertion of the lemma is not valid (exercise 14 of IV, p. 316).

PROPOSITION 22. --- Suppose that \(u_{\mathrm{N}}\) is a positive endomorphism of \(\mathrm{L}^2(\mathrm{X})\). Then the endomorphism \(u_{\mathrm{N}}\) has finite trace if and only if the function \(x \mapsto \mathrm{N}(x,x)\) is \(\mu\) -integrable. In this case, we have

\[
\mathrm{Tr} (u _ {\mathrm{N}}) = \int_ {\mathrm{X}} \mathrm{N} (x, x) d \mu (x).
\]

If \(u_{N}\) has finite trace, Theorem 2 implies that \(x \mapsto \mathrm{N}(x, x)\) is integrable on X and that its integral is the trace of \(u_{N}\) .

Conversely, suppose that the function \(x \mapsto \mathrm{N}(x, x)\) is integrable. Since \(u_{N}\) is positive, the orthonormal families \((f_i)_{i \in I}\) and \((g_i)_{i \in I}\) can be chosen such that \(f_i = g_i\) for all \(i \in I\) (remark 3 of IV, p. 151). For every \(x \in \widetilde{X}\) , we have

\[
\mathrm{N} (x, x) = \sum_ {i \in \mathrm{I}} \alpha_ {i} | f _ {i} (x) | ^ {2},
\]

whence, for every finite subset J of I, we deduce

\[
\begin{array}{c} \sum_ {i \in \mathrm{J}} \alpha_ {i} = \sum_ {i \in \mathrm{J}} \alpha_ {i} \int_ {\mathrm{X}} | f _ {i} (x) | ^ {2} d \mu (x) \\ = \int_ {\mathrm{X}} \sum_ {i \in \mathrm{J}} \alpha_ {i} | f _ {i} (x) | ^ {2} d \mu (x) \leqslant \int_ {\mathrm{X}} \mathrm{N} (x, x) d \mu (x). \end{array}
\]

The family \((\alpha_{i})_{i\in\mathrm{I}}\) is therefore summable, which implies that the endomorphism \(u_{N}\) has finite trace (prop. 12 of IV, p. 164).

Remark. --- Even when X is compact and N is continuous, the endomorphism \(u_{\mathrm{N}}\) of \(\mathrm{L}^{2}(\mathrm{X})\) is not always of finite trace (cf. exercise 8 of IV, p. 314).

